% GENERATED FILE. Edit canonical lessons, then run npm run generate:books.
% canonical-source-hash: fnv1a64:1df886c8001eb651
% canonical-lessons: PT-C314-falecimento, PT-C314-falecer, PT-C314-obito, PT-C314-luto, PT-C314-velorio

\chapter{Uma perda na família}
\label{ch:pt-uma-perda-na-familia}
\hlchaptermodality{314}

\begin{chapteropening}
\textbf{By the end of this chapter:} \emph{I can talk about a death in the family, mourning and the wake, in respectful words.}

Complete the last lesson of chapter 314: I can talk about a death in the family, mourning and the wake, in respectful words.
\end{chapteropening}

\section[o falecimento]{o falecimento — the passing, the death (formal)}

\label{lesson:PT-C314-falecimento}



\begin{quote}
Before the new one: say the Portuguese for the separation, then the Portuguese for the move.
\end{quote}

\subsection*{You'll want to know: o falecimento}
\textbf{o falecimento} — \textquotedblleft{}the passing, the death (formal)\textquotedblright{}.

\begin{quote}
O falecimento do avô juntou a família toda.
\end{quote}

\begin{quote}
\emph{Grandfather’s passing brought the whole family together.}
\end{quote}

A respectful word, common in notices.

\begin{cousinweb}
\emph{falecimento} is built on \emph{falecer} \textquotedblleft{}to pass away\textquotedblright{}, with the ending \emph{-mento}; the verb goes back to Vulgar Latin \emph{fallēscere} (EtymDB).
\end{cousinweb}

\begin{grammarlens}[title={What you've built}]
One more word for telling, explaining and describing at B1.
\end{grammarlens}

\subsection*{Your turn}
\begin{itemize}
\raggedright
  \item \emph{Say it:} \emph{o falecimento}
  \item \emph{Say it:} \emph{o falecimento}, once more
  \item \emph{Say it:} a sentence of your own with \emph{o falecimento}
  \item \emph{Recall:} say the Portuguese for the separation, then the Portuguese for the move, then say \emph{o falecimento} again
\end{itemize}

\begin{tcolorbox}[breakable,colback=teal!4,colframe=teal!35!black,title={Before you move on}]
What does o falecimento mean? (\textquotedblleft{}The passing, the death (formal)\textquotedblright{}.) Say it once more.
\end{tcolorbox}
\section[falecer]{falecer — to pass away}

\label{lesson:PT-C314-falecer}



\begin{quote}
Before the new one: say the Portuguese for the move, then the Portuguese for the passing.
\end{quote}

\subsection*{You'll want to know: falecer}
\textbf{falecer} — \textquotedblleft{}to pass away\textquotedblright{}.

\begin{quote}
O avô dela faleceu no domingo.
\end{quote}

\begin{quote}
\emph{Her grandfather passed away on Sunday.}
\end{quote}

A gentle, formal way to say \emph{morrer}.

\begin{cousinweb}
\emph{falecer} goes back to Vulgar Latin \emph{fallēscere} (EtymDB), from Latin \emph{fallere} \textquotedblleft{}to deceive, to fail\textquotedblright{}: life fails.
\end{cousinweb}

\begin{grammarlens}[title={What you've built}]
One more word for telling, explaining and describing at B1.
\end{grammarlens}

\subsection*{Your turn}
\begin{itemize}
\raggedright
  \item \emph{Say it:} \emph{falecer}
  \item \emph{Say it:} \emph{falecer}, once more
  \item \emph{Say it:} a sentence of your own with \emph{falecer}
  \item \emph{Recall:} say the Portuguese for the move, then the Portuguese for the passing, then say \emph{falecer} again
\end{itemize}

\begin{tcolorbox}[breakable,colback=teal!4,colframe=teal!35!black,title={Before you move on}]
What does falecer mean? (\textquotedblleft{}To pass away\textquotedblright{}.) Say it once more.
\end{tcolorbox}
\section[o óbito]{o óbito — the death (in official language)}

\label{lesson:PT-C314-obito}



\begin{quote}
Before the new one: say the Portuguese for the passing, then the Portuguese for to pass away.
\end{quote}

\subsection*{You'll want to know: o óbito}
\textbf{o óbito} — \textquotedblleft{}the death (in official language)\textquotedblright{}.

\begin{quote}
O médico confirmou o óbito às seis horas.
\end{quote}

\begin{quote}
\emph{The doctor confirmed the death at six o’clock.}
\end{quote}

An official word: \emph{certidão de óbito} — death certificate.

\begin{cousinweb}
\emph{óbito} is a learned borrowing from Latin \emph{obitus} \textquotedblleft{}a going down, death\textquotedblright{}, from \emph{obīre} \textquotedblleft{}to go towards, to die\textquotedblright{}, built on \emph{īre} \textquotedblleft{}to go\textquotedblright{} (Lewis \& Short).
\end{cousinweb}

\begin{grammarlens}[title={What you've built}]
One more word for telling, explaining and describing at B1.
\end{grammarlens}

\subsection*{Your turn}
\begin{itemize}
\raggedright
  \item \emph{Say it:} \emph{o óbito}
  \item \emph{Say it:} \emph{o óbito}, once more
  \item \emph{Say it:} a sentence of your own with \emph{o óbito}
  \item \emph{Recall:} say the Portuguese for the passing, then the Portuguese for to pass away, then say \emph{o óbito} again
\end{itemize}

\begin{tcolorbox}[breakable,colback=teal!4,colframe=teal!35!black,title={Before you move on}]
What does o óbito mean? (\textquotedblleft{}The death (in official language)\textquotedblright{}.) Say it once more.
\end{tcolorbox}
\section[o luto]{o luto — mourning}

\label{lesson:PT-C314-luto}



\begin{quote}
Before the new one: say the Portuguese for to pass away, then the Portuguese for the death.
\end{quote}

\subsection*{You'll want to know: o luto}
\textbf{o luto} — \textquotedblleft{}mourning\textquotedblright{}.

\begin{quote}
A família está de luto.
\end{quote}

\begin{quote}
\emph{The family is in mourning.}
\end{quote}

\emph{Estar de luto} — to be in mourning; people often wear black.

\begin{cousinweb}
\emph{luto} is inherited from Latin \emph{lūctus} \textquotedblleft{}grief, sorrow\textquotedblright{} (EtymDB).
\end{cousinweb}

\begin{grammarlens}[title={What you've built}]
One more word for telling, explaining and describing at B1.
\end{grammarlens}

\subsection*{Your turn}
\begin{itemize}
\raggedright
  \item \emph{Say it:} \emph{o luto}
  \item \emph{Say it:} \emph{o luto}, once more
  \item \emph{Say it:} a sentence of your own with \emph{o luto}
  \item \emph{Recall:} say the Portuguese for to pass away, then the Portuguese for the death, then say \emph{o luto} again
\end{itemize}

\begin{tcolorbox}[breakable,colback=teal!4,colframe=teal!35!black,title={Before you move on}]
What does o luto mean? (\textquotedblleft{}Mourning\textquotedblright{}.) Say it once more.
\end{tcolorbox}
\section[o velório]{o velório — the wake}

\label{lesson:PT-C314-velorio}



\begin{quote}
Before the new one: say the Portuguese for the death, then the Portuguese for mourning.
\end{quote}

\subsection*{You'll want to know: o velório}
\textbf{o velório} — \textquotedblleft{}the wake\textquotedblright{}.

\begin{quote}
Muitos vizinhos foram ao velório.
\end{quote}

\begin{quote}
\emph{Many neighbours went to the wake.}
\end{quote}

The night when people keep watch with the dead before the burial.

\begin{cousinweb}
\emph{velório} is built on \emph{velar} \textquotedblleft{}to keep watch\textquotedblright{} with the ending \emph{-ório}; \emph{velar} comes from Latin \emph{vigilāre} \textquotedblleft{}to stay awake, to keep watch\textquotedblright{} (Lewis \& Short).
\end{cousinweb}

\begin{grammarlens}[title={What you've built}]
One more word for telling, explaining and describing at B1.
\end{grammarlens}

\subsection*{Your turn}
\begin{itemize}
\raggedright
  \item \emph{Say it:} \emph{o velório}
  \item \emph{Say it:} \emph{o velório}, once more
  \item \emph{Say it:} a sentence of your own with \emph{o velório}
  \item \emph{Recall:} say the Portuguese for the death, then the Portuguese for mourning, then say \emph{o velório} again
\end{itemize}

\begin{tcolorbox}[breakable,colback=teal!4,colframe=teal!35!black,title={Before you move on}]
What does o velório mean? (\textquotedblleft{}The wake\textquotedblright{}.) Say it once more.
\end{tcolorbox}
